\section{Background}
\label{sec:background}

\begin{intuition}
This section collects the handful of ideas we lean on throughout: what
an embedding is, what it means to \emph{cluster} a memory, how we
measure how ``spread out'' a cluster is ($\dbar$), and what effective
dimension ($\deff$) captures about a cluster's shape. A reader
familiar with retrieval-augmented generation, product quantization,
and effective rank can skim. The rest of the paper assumes only these.
\end{intuition}

This section exists for a reader approaching the paper from an adjacent
field. A reader already familiar with retrieval-augmented generation,
product quantization, and effective rank can safely skim.

\subsection{What is an embedding, and what is an embedding memory?}
\label{sec:background:embedding}

An \emph{embedding} is a function that maps a piece of input data --- a
sentence, an image, a code snippet --- to a point in $\R^d$ for some
fixed $d$ between a few hundred and a few thousand. Modern sentence
embeddings have two properties that matter for this paper. First, they
are \emph{normalised}: the embedded vectors live on the unit sphere
$\SphereD$, so similarity reduces to the cosine $\inner{x}{y} = \cos
\angle(x, y)$. Second, they are \emph{semantically structured}: two
sentences that mean the same thing have cosine similarity close to one,
while two unrelated sentences have cosine similarity close to zero. The
second property is enforced during training by a contrastive
loss~\cite{reimers2019sbert,gao2021simcse,xiao2023bge,wang2022e5} that
pulls paraphrases together and pushes unrelated pairs apart.

An \emph{embedding memory} is any system that stores a large set of such
vectors $\memset = \{x_1, \ldots, x_N\} \subset \SphereD$ and supports
nearest-neighbour queries on them. The three canonical systems in the
introduction --- RAG databases, hippocampal memory, continual-learning
replay buffers --- all take this form, differing mainly in how they choose
what to keep.

\subsection{Retrieval-augmented generation, quickly}
\label{sec:background:rag}

Retrieval-augmented generation
(RAG)~\cite{lewisRAG2020,guu2020realm,karpukhin2020dpr,izacardAtlas2023,
borgeaud2022retro,izacard2021fid} is the dominant way large language
models access factual knowledge outside their training distribution. The
pipeline has three stages:

\begin{enumerate}[leftmargin=1.8em, itemsep=2pt]
\item \textbf{Indexing.} Every passage in the corpus (often the whole of
Wikipedia, or a domain-specific collection) is passed through a
sentence encoder and the resulting vectors are stored in a
vector database.
\item \textbf{Retrieval.} At query time, the query string is embedded
by the same encoder, and the $k$ passages whose vectors are most
similar to the query's vector are retrieved by approximate
nearest-neighbour search~\cite{johnson2019faiss,malkovHNSW2018}.
\item \textbf{Reading.} The retrieved passages are concatenated into
the context of a large language model, which generates the final
answer.
\end{enumerate}

At production scale the vector database has $10^{6}$--$10^{10}$
entries. Scanning all of them for every query is infeasible, so every
production stack compresses the index. The dominant compression family
is \emph{vector quantization}: approximate each $d$-dimensional vector
by a short code, so the database is a few bytes per vector rather than
$4d$. Two examples of this family dominate:

\paragraph{Product Quantization (PQ).} Introduced
in~\cite{jegouPQ2011}, PQ splits each vector into $M$ sub-vectors of
dimension $d/M$, and quantizes each sub-vector to one of $K$ centroids
learned by $k$-means. The code for each vector is $M \log_2 K$ bits ---
for example, $M=8$, $K=256$ gives $64$ bits per $d$-dimensional vector.
Distances are then computed as the sum of pre-tabulated distances
between sub-vector centroids. Optimized
PQ~\cite{ge2013optimized} adds a learned rotation to decorrelate
sub-vectors before chunking. These two remain the backbone of
billion-scale similarity search~\cite{johnson2019faiss}.

\paragraph{Locality-Sensitive Hashing (LSH).} LSH for cosine similarity
was introduced in~\cite{charikarLSH2002} and generalised
in~\cite{datarLSH2004}. The idea is to hash each vector to a short
binary string such that vectors with high cosine similarity collide
with high probability. Matching becomes bitwise Hamming-distance
comparison, which is cheap. LSH is simpler than PQ but sits below
PQ on the Pareto frontier for typical ANN
workloads~\cite{johnson2019faiss}.

\paragraph{Graph-based ANN.} HNSW~\cite{malkovHNSW2018} and its
descendants store not vectors but a navigable small-world graph whose
edges connect nearby vectors. Retrieval is a greedy walk from an entry
point. HNSW at full precision is typically the fastest recall-$95$
option, but its memory cost is dominated by the graph, so at large
scale it is combined with PQ or with HNSW-prune (a retention sweep that
drops the lowest-degree fraction).

All three families compress \emph{at the vector level}: every vector
stays in the index, only represented more cheaply. What this paper calls
\emph{consolidation} is orthogonal: it compresses \emph{at the cluster
level}, dropping most vectors and replacing them with a handful of
representatives. Whether consolidation should win over vector quantization
depends on the geometry of the corpus, and is the central empirical
question of Section~\ref{sec:e3}.

\subsection{Complementary learning systems, quickly}
\label{sec:background:cls}

The cognitive neuroscience literature on human memory consolidation is
anchored by the \emph{Complementary Learning Systems} (CLS) framework of
McClelland, McNaughton and
O'Reilly~\cite{mcclellandConsolidation1995}, itself drawing on Marr's
earlier model of the hippocampus as a fast associative
store~\cite{marrSimple1971} and on Alvarez and Squire's account of
gradient retrograde amnesia~\cite{alvarezSquireConsolidation1994}. The
claim: the brain has two memory systems. The hippocampus encodes
episodic traces rapidly but has limited capacity; the neocortex stores
overlapping, distributed representations and learns slowly. During
offline periods --- particularly sleep --- the hippocampus replays
compressed versions of its traces to the cortex, allowing slow cortical
consolidation without interference.

Key empirical anchors: Wilson and McNaughton showed that hippocampal
place cells replay awake-activity sequences during subsequent
sleep~\cite{wilson1994reactivation}. O'Reilly and
colleagues~\cite{oreillyHippocampus2014} laid out the computational
constraints on such a system --- sparse hippocampal coding, pattern
separation, pattern completion --- and argued that ``consolidation'' in
the neuroscience sense and ``rehearsal'' in the
continual-learning~\cite{mccloskey1989cf,robinsReplay1995} sense are the
same operation under different names.

Two open questions in the CLS literature are structurally similar to
the embedding consolidation problem we study. First, how many cortical
exemplars are needed to represent a hippocampal episode? Under common
assumptions this is set by task similarity and buffer size, but without
a geometric bound the number is chosen empirically. Second, what
happens when the episodes overlap heavily (e.g.\ thousands of
paraphrase-like traces of the same fact)? Standard CLS models assume
well-separated episodes consolidate without interference. We note that
our theorem, in the embedding-consolidation setting we study, gives a
geometric condition --- tightness of the cluster in its effective
subspace --- under which a structural analogue of that assumption
holds. Whether the same geometry governs biological consolidation is a
question the theorem \emph{motivates} but does not settle; the
evidence we report is for text embedding consolidation only.

\subsection{Effective dimension and why it matters}
\label{sec:background:deff}

The central spectral quantity of this paper is the \emph{effective
dimension} of a cluster, defined for $X \subset \R^d$ with sample
covariance $\Sigma$ and eigenvalues $\lambda_1 \geq \ldots \geq
\lambda_d \geq 0$ by
\begin{equation}
\deff(X) := \frac{\bigl(\sum_i \lambda_i\bigr)^2}{\sum_i \lambda_i^2}.
\label{eq:deff_def}
\end{equation}
This is the ratio between the squared trace of $\Sigma$ and the trace
of $\Sigma^2$; it lies in $[1, d]$ and takes the value $d$ iff $\Sigma$
is proportional to the identity, the value $1$ iff $\Sigma$ has a single
nonzero eigenvalue (rank-one collapse). In information-theoretic terms,
$1/\deff = \sum_i p_i^2$ where $p_i = \lambda_i / \sum_j \lambda_j$ is
the probability of eigenvalue $i$ under the natural Gibbs distribution;
$\deff$ is therefore the exponential of the R\'enyi-2 entropy of the
normalised spectrum~\cite{royVetterli2007}. In the
statistics literature the same quantity also goes by ``effective
degrees of freedom,''
``stable rank,'' or ``participation
ratio''~\cite{vershynin2018hdp,ledoux2001concentration}.

Two things make $\deff$ the right spectral summary for our problem.
\begin{enumerate}[leftmargin=1.8em, itemsep=2pt]
\item It is \emph{scale-invariant}: multiplying $\Sigma$ by a constant
leaves $\deff$ unchanged. The quantity therefore reports the shape of
the spectrum, not its magnitude.
\item It is \emph{weighted}: two large eigenvalues and one thousand
small ones give $\deff \approx 2$, not $\deff \approx 1000$. This
matches the geometric intuition that the ``hard to compress''
directions are those with large variance.
\end{enumerate}

The previous paper in this trilogy documented that production sentence
encoders have $\deff \approx 16$ at the global
level~\cite{geometryOfForgetting2026}, with $\deff_{\mathrm{local}}$ on a
single cluster often below $5$. The claim of the present paper is that
this same $\deff_{\mathrm{local}}$ controls how much compression a
cluster can take before its members stop retrieving themselves.

\paragraph{Measurement conventions (used throughout).}
All reported values of $\deff$ use Equation~\eqref{eq:deff_def}
(participation ratio); the code computes this form as the default of
\texttt{gac.theory.d\_eff} (see the repository). We distinguish:
(i) the \emph{global} $\deff$ of a corpus is measured on the pooled
embedding matrix of all items across clusters;
(ii) the \emph{local} $\deff_{\mathrm{local}}$ of a corpus is the
\textbf{median} of per-cluster participation ratios, computed on each
cluster's own embedding matrix.
The theorem as stated is cluster-conditional: its bound applies to a
single cluster with $\deff$ computed from that cluster's covariance.
Corpus-level diagnostics (e.g.\ the ordering tables in
Section~\ref{sec:template_collapse}) report medians across clusters.
Likewise $\dbar$ is mean within-cluster cosine distance and is reported
cluster-conditionally; corpus rows report medians. All other
definitions used in this paper reduce to this single quantity $\deff$
under the participation-ratio form; we do not mix definitions.

\subsection{Concentration of measure on the sphere}
\label{sec:background:concentration}

The proof of the duality theorem uses a classical fact about points on a
high-dimensional sphere: most of them are nearly orthogonal to any fixed
direction. Formally, if $x$ is drawn uniformly from $\SphereD$, then for
any fixed $r \in \SphereD$ and any $\theta \in (0, 1)$,
\begin{equation}
\Pr[\inner{x}{r} \geq \theta] \;\leq\; e^{-d \theta^2 / 2}.
\label{eq:cap_volume_standard}
\end{equation}
This is the canonical \emph{cap-volume} estimate (see Ball's
introduction to convex geometry~\cite{ball1997elementary} or Vershynin's
textbook~\cite{vershynin2018hdp}; tight constants appear in Milman and
Schechtman~\cite{milman1986asymptotic} and in
Ledoux~\cite{ledoux2001concentration}).

Equation~\eqref{eq:cap_volume_standard} says: on a high-dimensional
sphere, the spherical cap of angular radius $\arccos\theta$ has
exponentially small mass as $d$ grows. This is why nearest-neighbour
search on random points in high dimensions becomes
trivial~\cite{beyer1999nn}: no point is close to the query except the
query itself, if there is one.

Our setting is different in one crucial way. Our cluster is not isotropic
on $\SphereD$; it concentrates in a low-dimensional effective subspace of
dimension $\deff \ll d$. The bound we need therefore replaces $d$ by
$\deff$: on a cluster with effective dimension $\deff$ and mean
within-cluster distance $\dbar$, a spherical cap of radius
$\arccos\theta$ has mass at most $c_1 (\thetap/\dbar)^{\deff/2}$
whenever $\thetap < \dbar$ (Lemma~\ref{lem:cap} in Methods). A union
bound over $m$ representatives then gives the theorem. The
low-dimensionality of real embedding clusters is what makes the bound
non-trivial: for $\deff = 16$ and $\thetap / \dbar = 0.5$, the
per-representative mass bound is $(0.5)^{8} = 0.004$ --- small, but not
astronomically small. For $\deff = 2$ (HotpotQA-like), $(0.5)^1 = 0.5$,
and the bound is vacuous within a factor of two.

\subsection{Notation summary}
\label{sec:background:notation}

We use the following notation throughout.

\begin{center}
\begin{tabular}{ll}
\toprule
Symbol & Meaning \\
\midrule
$\clusterset = \{x_1, \ldots, x_n\} \subset \SphereD$ & cluster of $n$
source points on the unit sphere \\
$\reps = \{r_1, \ldots, r_m\} \subset \SphereD$ & set of $m$
representatives produced by a consolidator \\
$\Sigma$ & sample covariance of $\clusterset$ \\
$\lambda_1 \geq \ldots \geq \lambda_d$ & eigenvalues of $\Sigma$ \\
$\deff = (\sum_i \lambda_i)^2 / \sum_i \lambda_i^2$ & effective dimension
of $\clusterset$ (Def.~\ref{def:deff}) \\
$\dbar = \tfrac{1}{\binom{n}{2}} \sum_{i<j} (1 - \inner{x_i}{x_j})$ &
mean within-cluster cosine distance \\
$\theta \in (0, 1)$ & retrieval cap (minimum cosine to count as a hit) \\
$\thetap = 1 - \theta$ & cap slack \\
$\eps_{\mathrm{cap}}$ & cap-coverage error (Eq.~\ref{eq:cap_error}) \\
$\eps_{\mathrm{id}}$ & identity-retrieval error (Def.~\ref{def:id}) \\
$\rho = \lambda_1 / \sum_i \lambda_i$ & spectral concentration (top eigenvalue fraction) \\
\bottomrule
\end{tabular}
\end{center}
